\section*{अध्याय \ref{ch:b2:equilibrium-shifts} --- साम्य स्थिरांक और साम्य का विस्थापन}

\begin{solution}{exo:b2:equilibrium-shifts:1}
$K^\circ = \exp(32\,800/(8.314 \times 298.15)) = \num{5.6e5}$, और
$+\qty{32.8}{kJ/mol}$ के लिए इसका व्युत्क्रम, \num{1.8e-6}। 100 के गुणक के लिए
$\Delta(\Delta_r G^\circ) = -RT\ln 100 = \qty{-11.4}{kJ/mol}$ चाहिए।
\end{solution}

\begin{solution}{exo:b2:equilibrium-shifts:2}
$\ln K^\circ(400) = \ln(5.4 \times 10^5) + (91\,880/8.314)(1/400 -
1/298.15) = 13.20 - 9.44 = 3.76$, $K^\circ \approx 43$। सारणियाँ 35.6 देती हैं:
\qty{100}{K} में ही स्थिर-$\Delta_r H^\circ$ सन्निकटन 20\,\% हट जाता है
(ताप $T$ बढ़ने पर अभिक्रिया अधिक \omterm{def:b2:reaction-enthalpy:exothermic}{ऊष्माक्षेपी} होती जाती है)।
\end{solution}

\begin{solution}{exo:b2:equilibrium-shifts:3}
(a) $n = 1$, $r = 0$, $\varphi = 2$: $v = 1$ (वाष्प दाब
$T$ का फलन है)। (b) $v = 3 - 1 + 2 - 3 = 1$। (c) $v = 3 - 1 + 2 - 1 = 3$। (d) चार
स्पीशीज़, एक अभिक्रिया, दो प्रावस्थाएँ (कार्बन और गैस): $v = 4 - 1 + 2 - 2 =
3$; यदि गैस केवल कार्बन और भाप से बनी है, तो गैस में $n(\ce{CO}) = n(\ce{H2})$,
$s = 1$ और $v = 2$।
\end{solution}

\begin{solution}{exo:b2:equilibrium-shifts:4}
गर्म करना: \omterm{def:b2:reaction-enthalpy:exothermic}{ऊष्माशोषी} दिशा, \ce{NO2} की ओर (अग्र)। संपीडन:
कम गैस अणुओं की ओर, \ce{N2O4} (पश्च)। नियत आयतन पर आर्गन:
कोई परिवर्तन नहीं, अभिक्रियारत गैसों के आंशिक दाब अपरिवर्तित रहते हैं।
\end{solution}

\begin{solution}{exo:b2:equilibrium-shifts:5}
\qty{298}{K} पर: $\Delta_r G^\circ = 180.6 - 298.15 \times 0.0248 =
\qty{173.2}{kJ/mol}$, $K^\circ = \num{4.5e-31}$। \qty{2000}{K} पर: $\Delta_r
G^\circ = 180.6 - 2000 \times 0.0248 = \qty{131.0}{kJ/mol}$, $K^\circ =
\num{3.8e-4}$। $\Delta\nu_{\text{गैस}} = 0$ के साथ $x(\ce{NO})^2 = K^\circ x(\ce{N2})
x(\ce{O2})$: $x(\ce{NO}) = \sqrt{3.8 \times 10^{-4} \times 0.78 \times 0.21} =
\num{7.9e-3}$, लगभग एक प्रतिशत। इंजन ऐसे तापों पर जलाते हैं; गैसों के
ठंडा होने पर $K^\circ$ बहुत अधिक गिरता है, परंतु \ce{NO} का विघटन
अत्यंत मंद है: गर्म साम्य निकास में जमा हुआ रह जाता है (जब तक कोई उत्प्रेरक
उसे हटा न दे)।
\end{solution}

\begin{solution}{exo:b2:equilibrium-shifts:6}
एस्टर के $\xi$ mol और स्थिर कुल मात्रा के साथ $Q =
\xi^2/[(1 - \xi)(n_B - \xi)]$। $n_B = 1.0$ के लिए: $\xi^2/(1 - \xi)^2 = 4$, $\xi
= \qty{0.67}{mol}$। $n_B = 3.0$ के लिए: $3\xi^2 - 16\xi + 12 = 0$, $\xi =
\qty{0.90}{mol}$। जल हटाना $Q$ को $K^\circ$ से नीचे रखता है, $\xi$ चाहे जो हो:
अभिक्रिया अम्ल के उपभुक्त होने तक अग्र दिशा में चलती रहती है।
\end{solution}

\begin{solution}{exo:b2:equilibrium-shifts:7}
आर्गन के बिना: मात्राएँ $1 - \alpha$, $2\alpha$, कुल $1 + \alpha$;
$Q = 4\alpha^2/(1 - \alpha^2) = 0.148$, $\alpha = \sqrt{0.148/4.148} = 0.189$।
\qty{1}{bar} पर \qty{1}{mol} आर्गन के साथ: कुल $2 + \alpha$, $Q =
4\alpha^2/[(1 - \alpha)(2 + \alpha)] = 0.148$, अतः $4.148\alpha^2 + 0.148\alpha
- 0.296 = 0$, $\alpha = 0.250$: नियत कुल दाब पर तनु करना
आंशिक दाबों को घटाता है, जो अधिक गैस अणुओं वाले पक्ष के अनुकूल है। नियत
आयतन पर अभिक्रियारत गैसों के आंशिक दाब नहीं बदलते:
$\alpha$ 0.189 ही रहता है।
\end{solution}

\begin{solution}{exo:b2:equilibrium-shifts:8}
केवल ठोस: $n = 3$, $r = 1$, $s = 1$ ($p(\ce{NH3}) = p(\ce{HCl})$),
$\varphi = 2$: $v = 3 - 1 - 1 + 2 - 2 = 1$; $T$ चुनिए और दाब
नियत हो जाता है (एक ``वियोजन दाब'')। अमोनिया मिलाने पर $s = 0$ और $v = 2$:
$T$ और एक दाब चुना जा सकता है; गुणनफल $p(\ce{NH3})\,p(\ce{HCl})$
$T$ द्वारा नियत होता है।
\end{solution}

\begin{solution}{exo:b2:equilibrium-shifts:9}
$\Delta_r H^\circ = -R\ln(K_2/K_1)/(1/T_2 - 1/T_1) = -8.314 \times
\ln(0.0173)/(1/600 - 1/500) = \qty{-101}{kJ/mol}$, जो
\qty{298}{K} वाले मान (\qty{-91.9}{kJ/mol}) से अधिक ऋणात्मक है, जैसा \omterm{thm:b2:reaction-enthalpy:kirchhoff}{किरखॉफ़ के नियम} ने पूर्वानुमान दिया था
($\Delta_r C_p^\circ < 0$)।
\end{solution}

\begin{solution}{exo:b2:equilibrium-shifts:10}
$Q = \dfrac{n_{\ce{NH3}}^2\,n_{\text{कुल}}^2}{n_{\ce{N2}}n_{\ce{H2}}^3}
\left(\dfrac{p^\circ}{p}\right)^2$। नियत $T$, $p$ पर:
$\partial\ln Q/\partial n_{\ce{N2}} = -1/n_{\ce{N2}} + 2/n_{\text{कुल}}$, जो धनात्मक है
जब $n_{\ce{N2}} > n_{\text{कुल}}/2$, अर्थात् $x(\ce{N2}) > 1/2$। तब $Q$,
$K^\circ$ से अधिक हो जाता है और साम्य पश्च दिशा में विस्थापित होता है। नियत आयतन पर
$\partial\ln Q/\partial n_{\ce{N2}} = -1/n_{\ce{N2}} < 0$: सदा अग्र।
\end{solution}

\begin{solution}{exo:b2:equilibrium-shifts:11}
$K^\circ = \exp(-1620/(8.314 \times 1100)) = 0.84$: साम्य पर $p(\ce{CO2}) =
\qty{0.84}{bar}$, अर्थात् गैस के $n = pV/RT = 0.84 \times 10^5 \times
0.0100/(8.314 \times 1100) = \qty{0.092}{mol}$। चूना पत्थर के \qty{0.050}{mol}
के साथ सारा विघटित हो जाता है ($p = \qty{0.46}{bar} < K^\circ p^\circ$):
कोई साम्य नहीं, साम्य का भंग। \qty{0.200}{mol} के साथ साम्य:
\ce{CO2} और \ce{CaO} के \qty{0.092}{mol}, \ce{CaCO3} के \qty{0.108}{mol} शेष।
\end{solution}

\begin{solution}{exo:b2:equilibrium-shifts:12}
संस्तर 1 लगभग \qty{890}{K} और 68\,\% रूपांतरण पर निकलता है, संस्तर 2 लगभग
\qty{785}{K} और 92\,\% पर, संस्तर 3 लगभग \qty{725}{K} और 98\,\% पर। एक
रुद्धोष्म संस्तर में गैस अभिक्रिया करते हुए तब तक गर्म होती है जब तक वह साम्य
वक्र से न मिले, जो उस ताप पर केवल लगभग 68\,\% की अनुमति देता है। पूरे समय ठंडा
उत्प्रेरक कोई विकल्प नहीं, क्योंकि उत्प्रेरक केवल गर्म होने पर काम करता है; संस्तरों
के बीच ठंडा करना दर को बनाए रखता है और साम्य की सीमा को ऊपर ले जाता है। अंतिम संस्तर
उस निम्नतम ताप के साम्य से सीमित है जिस पर उत्प्रेरक
अब भी काम करता है (संयंत्र तब ट्राइऑक्साइड को अवशोषित करके गैस को एक बार फिर
उत्प्रेरक पर से गुज़ारते हैं)।
\end{solution}

\begin{solution}{pb:b2:equilibrium-shifts:1}
\textbf{1.} $\Delta_r H^\circ = \qty{-91.88}{kJ/mol}$, $\Delta_r S^\circ =
2(192.77) - 191.61 - 3(130.68) = \qty{-198.1}{J/(K.mol)}$, $\Delta_r G^\circ =
\qty{-32.8}{kJ/mol}$, $K^\circ = \num{5.6e5}$।
\textbf{2.} $K^\circ(700) = \exp(-54\,380/(8.314 \times 700)) = \num{8.8e-5}$;
$K^\circ(800) = \exp(-77\,320/(8.314 \times 800)) = \num{8.9e-6}$।
\textbf{3.} $\Delta_r H^\circ = -R\ln(K_{800}/K_{700})/(1/800 - 1/700)$,
अर्थात् \qty{-106}{kJ/mol}, जो \qty{298}{K} वाले मान से अधिक ऋणात्मक है।
\textbf{4.} $\ln K^\circ(723) = \ln(8.75 \times 10^{-5}) + 12\,777 \times
(1/723.15 - 1/700)$, जहाँ $12\,777 = 106\,230/8.314$: $K^\circ =
\num{4.9e-5}$।
\textbf{5.} $\Delta_r G^\circ(723) = -91.88 + 723.15 \times 0.1981 =
\qty{51.4}{kJ/mol}$, $K^\circ = \num{1.9e-4}$: चार गुना बहुत बड़ा, क्योंकि
\qty{425}{K} में $\Delta_r H^\circ$ और $\Delta_r S^\circ$ बदलते हैं
($\Delta_r C_p^\circ \approx \qty{-44}{J/(K.mol)}$)।
\textbf{6.} \ce{N2} $1 - \xi$, \ce{H2} $3 - 3\xi$, \ce{NH3} $2\xi$, कुल $4 -
2\xi$।
\textbf{7.} $x = 2\xi/(4 - 2\xi)$, अतः $1 - x = (4 - 4\xi)/(4 - 2\xi)$;
नाइट्रोजन का अंश $(1 - \xi)/(4 - 2\xi) = (1 - x)/4$ और हाइड्रोजन का अंश
उसका तिगुना है।
\textbf{8.}
\[
  K^\circ = \frac{x^2}{\frac{1-x}{4}\left(\frac{3(1-x)}{4}\right)^3}
  \left(\frac{p^\circ}{p}\right)^2 = \frac{256}{27}\,\frac{x^2}{(1 - x)^4}
  \left(\frac{p^\circ}{p}\right)^2 ;
\]
वर्गमूल लीजिए।
\textbf{9.} $a = (\sqrt{27}/16) \times \sqrt{4.9 \times 10^{-5}} \times 1 =
\num{2.27e-3}$; $x/(1 - x)^2 = a$ से $x = 0.0023$।
\textbf{10.} $a = 0.455$; $ax^2 - (2a + 1)x + a = 0$: $x = 0.25$।
\textbf{11.} $n = 3$, $r = 1$, $s = 1$ (1 : 3 का अनुपात अभिक्रिया द्वारा बना रहता है),
$\varphi = 1$: $v = 2$। $T$ और $p$ चुन लिए जाने पर कुछ भी स्वतंत्र नहीं बचता।
\textbf{12.} दाब: अग्र ($\Delta\nu_{\text{गैस}} = -2$)। ताप:
पश्च (\omterm{def:b2:reaction-enthalpy:exothermic}{ऊष्माक्षेपी})।
\textbf{13.} नियत कुल दाब पर अक्रिय गैस अभिक्रियारत गैसों के आंशिक
दाबों को घटाती है, जैसा दाब में गिरावट करती: पश्च।
\textbf{14.} $\partial\ln Q/\partial n_{\ce{N2}} = (-1/0.30 + 2)/n_{\text{कुल}} < 0$:
$Q$, $K^\circ$ से नीचे गिरता है, अग्र।
\textbf{15.} नहीं: नियत आयतन पर $\partial\ln Q/\partial n_{\ce{N2}} =
-1/n_{\ce{N2}} < 0$, संघटन चाहे जो हो: अग्र।
\textbf{16.} लगभग शून्य अमोनिया के साथ प्रवेश पर $Q$, $K^\circ$ से बहुत नीचे है:
गैस साम्य से बहुत दूर प्रवेश करती है और अग्र दिशा में अभिक्रिया करती है।
\textbf{17.} गैस उत्प्रेरक पर इतना कम समय बिताती है कि
साम्य नहीं पहुँचता; लंबा संपर्क एक बड़ा, अधिक महँगा परिवर्तक माँगता,
और लाभ थोड़ा होता।
\textbf{18.} $x = 2\xi/(4 - 2\xi) = 0.18$ से प्रति मोल फ़ीड \ce{N2}
$\xi = \qty{0.305}{mol}$: हाइड्रोजन (और नाइट्रोजन) का \omterm{def:b2:equilibrium-shifts:single-pass}{एकल-पारण रूपांतरण}
$0.305$, 31\,\%।
\textbf{19.} स्थायी संचालन में सारा ताज़ा फ़ीड अंततः रूपांतरित हो जाता है:
प्रति इकाई समय अमोनिया के \qty{2.00}{mol} निकलते हैं; परिवर्तक से प्रति इकाई समय \ce{N2} के $1.00/
0.305 = \qty{3.3}{mol}$ गुज़रने चाहिए, शेष पुनश्चक्रित होता है।
\textbf{20.} फ़ीड के साथ आने वाली आर्गन (और मेथेन) कभी अभिक्रिया नहीं करतीं और
द्रव अमोनिया के साथ नहीं हटतीं: \omterm{def:b2:equilibrium-shifts:single-pass}{निष्कासन} के बिना वे
पाश में संचित होकर अभिकारकों के आंशिक दाब घटा देतीं।
\textbf{21.} \qty{500}{K} पर लोहे का उत्प्रेरक बहुत मंद है: साम्य
अनुकूल तो होता, परंतु औद्योगिक समय में उसके निकट कभी नहीं पहुँचा जाता।
\textbf{22.} \qty{723}{K} और \qty{200}{bar} पर, रससमीकरणमितीय फ़ीड, आदर्श
गैसें: $\boldsymbol{x(\ce{NH3}) \approx 0.25}$।
\end{solution}
